% =====================================================================
\chapter{Composición e inversa}
\label{tema:4}
% =====================================================================

Componer funciones es aplicarlas una tras otra; invertir una función es deshacerla. Este tema estudia ambas operaciones y cómo se comportan respecto de la inyectividad y la sobreyectividad. Sigue las secciones~12.4 y~12.5 de Hammack~\cite{hammack}, las secciones~5.2 y~5.3 de Velleman~\cite{velleman} y la sección~4.3 de~\cite{mcs}.

\section{Composición}

\begin{definicion}[Composición]
Sean $f\colon A\to B$ y $g\colon B\to C$. La \emph{composición} de $g$ con $f$ es la función $g\circ f\colon A\to C$ definida por $(g\circ f)(a)=g(f(a))$.
\end{definicion}

El orden de lectura es de derecha a izquierda: en $g\circ f$ se aplica primero $f$. La composición solo está definida cuando el codominio de $f$ está contenido en el dominio de $g$.

\begin{ejemplo}\label{ej:t4-comp}
Sean $f,g\colon\R\to\R$, $f(x)=x^{2}$ y $g(x)=x+1$. Entonces $(g\circ f)(x)=x^{2}+1$ y $(f\circ g)(x)=(x+1)^{2}$. En $x=1$ valen $2$ y $4$: la composición \emph{no} es conmutativa.
\end{ejemplo}

\begin{proposicion}\label{prop:t4-asoc}
Sean $f\colon A\to B$, $g\colon B\to C$ y $h\colon C\to D$. Entonces $h\circ(g\circ f)=(h\circ g)\circ f$, y $f\circ\id_A=f=\id_B\circ f$.
\end{proposicion}

\begin{proof}
Ambas funciones tienen dominio $A$ y codominio $D$. Para $a\in A$, $\bigl(h\circ(g\circ f)\bigr)(a)=h\bigl((g\circ f)(a)\bigr)=h(g(f(a)))$ y $\bigl((h\circ g)\circ f\bigr)(a)=(h\circ g)(f(a))=h(g(f(a)))$. Las igualdades con la identidad son inmediatas.
\end{proof}

\section{Composición, inyectividad y sobreyectividad}

\begin{teorema}\label{teo:t4-comp}
Sean $f\colon A\to B$ y $g\colon B\to C$.
\begin{enumerate}
  \item Si $f$ y $g$ son inyectivas, $g\circ f$ es inyectiva.
  \item Si $f$ y $g$ son sobreyectivas, $g\circ f$ es sobreyectiva.
  \item Si $f$ y $g$ son biyectivas, $g\circ f$ es biyectiva.
\end{enumerate}
\end{teorema}

\begin{proof}
(i) Sean $a,a'\in A$ con $g(f(a))=g(f(a'))$. Como $g$ es inyectiva, $f(a)=f(a')$; como $f$ es inyectiva, $a=a'$. Los apartados (ii) y~(iii) son el \cref{eje:4-comp-sobre}.
\end{proof}

\begin{proposicion}\label{prop:t4-reciproco}
Sean $f\colon A\to B$ y $g\colon B\to C$.
\begin{enumerate}
  \item Si $g\circ f$ es inyectiva, $f$ es inyectiva.
  \item Si $g\circ f$ es sobreyectiva, $g$ es sobreyectiva.
\end{enumerate}
\end{proposicion}

\begin{proof}
(i) Si $f(a)=f(a')$, aplicando $g$ se obtiene $(g\circ f)(a)=(g\circ f)(a')$ y, por la inyectividad de $g\circ f$, $a=a'$. (ii) Sea $c\in C$. Por hipótesis existe $a\in A$ con $g(f(a))=c$; entonces $b=f(a)\in B$ cumple $g(b)=c$.
\end{proof}

Las otras dos implicaciones son falsas: si $g\circ f$ es inyectiva, $g$ no tiene por qué serlo, y si $g\circ f$ es sobreyectiva, $f$ no tiene por qué serlo (\cref{eje:4-cambridge}).

\section{Función inversa}

\begin{definicion}[Inversa]
Sea $f\colon A\to B$. Una función $g\colon B\to A$ es \emph{inversa} de $f$ si $g\circ f=\id_A$ y $f\circ g=\id_B$. Si solo se cumple $g\circ f=\id_A$, se dice que $g$ es una \emph{inversa por la izquierda} de $f$; si solo se cumple $f\circ g=\id_B$, una \emph{inversa por la derecha}.
\end{definicion}

\begin{teorema}\label{teo:t4-inversa}
Una función $f\colon A\to B$ tiene inversa si y solo si es biyectiva. En ese caso la inversa es única; se denota $f^{-1}$ y está dada por $f^{-1}(b)=$ el único $a\in A$ con $f(a)=b$.
\end{teorema}

\begin{proof}
($\Rightarrow$) Sea $g$ una inversa. Como $g\circ f=\id_A$ es inyectiva, $f$ es inyectiva por la \cref{prop:t4-reciproco}(i); como $f\circ g=\id_B$ es sobreyectiva, $f$ es sobreyectiva por la \cref{prop:t4-reciproco}(ii).

($\Leftarrow$) Si $f$ es biyectiva, cada $b\in B$ tiene exactamente una preimagen; sea $g(b)$ esa preimagen. Por construcción $f(g(b))=b$ para todo $b$. Además, para $a\in A$, $g(f(a))$ es la única preimagen de $f(a)$, que es $a$; luego $g\circ f=\id_A$.

(Unicidad) Si $g$ y $h$ son inversas de $f$, por la \cref{prop:t4-asoc}: $g=g\circ\id_B=g\circ(f\circ h)=(g\circ f)\circ h=\id_A\circ h=h$.
\end{proof}

\begin{proposicion}\label{prop:t4-laterales}
Sea $f\colon A\to B$ con $A\neq\emptyset$. Entonces $f$ tiene una inversa por la izquierda si y solo si es inyectiva, y tiene una inversa por la derecha si y solo si es sobreyectiva.
\end{proposicion}

\begin{proof}
Se demuestra la primera equivalencia; la segunda es el \cref{eje:4-derecha}. Si $g\circ f=\id_A$, $f$ es inyectiva por la \cref{prop:t4-reciproco}(i). Recíprocamente, sea $f$ inyectiva y fíjese $a_0\in A$. Se define $g\colon B\to A$ así: si $b\in f(A)$, $g(b)$ es la única preimagen de $b$; si $b\notin f(A)$, $g(b)=a_0$. Entonces $g(f(a))=a$ para todo $a\in A$.
\end{proof}

\begin{proposicion}\label{prop:t4-inv-comp}
Si $f\colon A\to B$ es biyectiva, $f^{-1}$ es biyectiva y $(f^{-1})^{-1}=f$. Si además $g\colon B\to C$ es biyectiva, $(g\circ f)^{-1}=f^{-1}\circ g^{-1}$.
\end{proposicion}

\begin{proof}
Las igualdades $f^{-1}\circ f=\id_A$ y $f\circ f^{-1}=\id_B$ dicen que $f$ es inversa de $f^{-1}$; por el \cref{teo:t4-inversa}, $f^{-1}$ es biyectiva y su inversa es $f$. La segunda afirmación es el \cref{eje:4-inv-comp}.
\end{proof}

\section{Cálculo de inversas}

Para invertir una biyección real $y=f(x)$ se despeja $x$ en función de $y$ y se comprueba que el resultado pertenece al dominio. Cuando $f$ no es biyectiva, puede \emph{restringirse} el dominio y el codominio para obtener una biyección: así se definen $\sqrt{\phantom{x}}$ como inversa de $x^{2}$ en $[0,\infty)$, $\ln$ como inversa de $\exp\colon\R\to(0,\infty)$ y $\operatorname{arcsen}$ como inversa de $\operatorname{sen}$ restringido a $[-\pi/2,\pi/2]$.

\begin{ejemplo}\label{ej:t4-inversa}
Sea $f\colon\R\to(1,\infty)$, $f(x)=1+e^{2x}$. Es biyectiva: de $y=1+e^{2x}$ con $y>1$ se obtiene $e^{2x}=y-1>0$ y $x=\tfrac12\ln(y-1)$. Luego $f^{-1}(y)=\tfrac12\ln(y-1)$. Comprobación: $f(f^{-1}(y))=1+e^{\ln(y-1)}=y$ y $f^{-1}(f(x))=\tfrac12\ln(e^{2x})=x$.
\end{ejemplo}

\begin{observacion}[Conexión con el curso]
Un tokenizador es un par $(\tau,\delta)$ con $\delta\circ\tau=\id_{\Sigma^{*}}$ (módulo~1, §1.2): la decodificación $\delta$ es una inversa por la izquierda de la codificación $\tau$. Por la \cref{prop:t4-reciproco}, $\tau$ es inyectiva y $\delta$ es sobreyectiva, sin que ninguna de las dos tenga que ser biyectiva. Una red neuronal profunda es una composición $f_L\circ\dots\circ f_2\circ f_1$ de capas (módulo~9), y la regla de la cadena para derivarla sigue esa misma estructura (módulos~4 y~10). La \cref{prop:t4-asoc} garantiza que el orden de agrupación de las capas es irrelevante.
\end{observacion}

\begin{observacion}[Errores frecuentes]
(1) Confundir $f^{-1}(x)$ con $1/f(x)$: para $f(x)=2x$, $f^{-1}(x)=x/2$ mientras que $1/f(x)=1/(2x)$. (2) Invertir el orden en la inversa de una composición: $(g\circ f)^{-1}=f^{-1}\circ g^{-1}$, no $g^{-1}\circ f^{-1}$ (para deshacer «ponerse los calcetines y luego los zapatos» se quitan primero los zapatos). (3) Escribir $f^{-1}$ para una función que no es biyectiva, en lugar de restringir antes su dominio o su codominio. (4) Creer que $g\circ f=\id_A$ basta para que $g$ y $f$ sean inversas: también hace falta $f\circ g=\id_B$.
\end{observacion}

\begin{lecturas}
\item Hammack~\cite{hammack}, §12.4 «Composition» y §12.5 «Inverse Functions».
\item Velleman~\cite{velleman}, §5.2 «One-to-One and Onto» y §5.3 «Inverses of Functions».
\item Lehman, Leighton y Meyer~\cite{mcs}, §4.3 «Functions».
\end{lecturas}

\section*{Ejercicios del tema 4}
\addcontentsline{toc}{section}{Ejercicios del tema 4}

\begin{ejercicio}{1}{eje:4-calc-comp}
Sean $f,g\colon\R\to\R$, $f(x)=x^{2}+1$ y $g(x)=2x-3$. Calcule $(g\circ f)(x)$, $(f\circ g)(x)$, $(f\circ f)(x)$ y $(g\circ g)(x)$, y evalúe las cuatro en $x=1$.
\end{ejercicio}

\begin{ejercicio}{1}{eje:4-vf}
Decida si son verdaderas o falsas: (a) $f\circ g=g\circ f$ para todas las funciones $f,g\colon\R\to\R$; (b) si $f\colon A\to B$ es biyectiva, $f^{-1}\circ f=f\circ f^{-1}$; (c) $f^{-1}(x)=1/f(x)$ para toda biyección $f\colon(0,\infty)\to(0,\infty)$; (d) $(f\circ g)\circ h=f\circ(g\circ h)$ siempre que las composiciones estén definidas.
\end{ejercicio}

\begin{ejercicio}{2}{eje:4-tablas}
Sean $f,g\colon\{1,2,3\}\to\{1,2,3\}$ con $f(1)=2$, $f(2)=3$, $f(3)=1$ y $g(1)=1$, $g(2)=1$, $g(3)=3$. Calcule $g\circ f$, $f\circ g$, $f\circ f$ y $f\circ f\circ f$. ¿Cuáles de $f$, $g$, $g\circ f$ tienen inversa? Calcule $f^{-1}$.
\end{ejercicio}

\begin{ejercicio}{2}{eje:4-inversas-reales}
Demuestre que son biyectivas y calcule sus inversas, comprobando las dos composiciones: (a) $f\colon\R\to\R$, $f(x)=2x-7$; (b) $g\colon(0,\infty)\to\R$, $g(x)=\ln(3x)$; (c) $h\colon[1,\infty)\to[0,\infty)$, $h(x)=\sqrt{x-1}$.
\end{ejercicio}

\begin{ejercicio}{2}{eje:4-inversa-z2}
Calcule la inversa de la biyección $f\colon\Z\times\Z\to\Z\times\Z$, $f(m,n)=(m+n,\ m+2n)$, del \cref{eje:3-lineal-z2}, y compruebe que $f^{-1}\circ f=\id$ y $f\circ f^{-1}=\id$.
\end{ejercicio}

\begin{ejercicio}{3}{eje:4-comp-sobre}
Sean $f\colon A\to B$ y $g\colon B\to C$. Demuestre que si $f$ y $g$ son sobreyectivas, entonces $g\circ f$ es sobreyectiva, y deduzca que la composición de dos biyecciones es una biyección.
\fuente{adaptado de Hammack~\cite{hammack}, §12.4 «Composition», ejercicios de la sección.}
\end{ejercicio}

\begin{ejercicio}{3}{eje:4-cambridge}
Sean $g\colon A\to B$ y $f\colon B\to C$. (a) Si $f\circ g$ es inyectiva, ¿debe ser $f$ inyectiva? ¿Y $g$? (b) Si $f\circ g$ es sobreyectiva, ¿debe ser $f$ sobreyectiva? ¿Y $g$? Demuestre las afirmaciones ciertas y dé contraejemplos de las falsas.
\fuente{adaptado de Cambridge, IA Numbers and Sets~\cite{camns}, Example Sheet~1, pregunta~8.}
\end{ejercicio}

\begin{ejercicio}{3}{eje:4-iteradas}
Para $f\colon A\to A$ se escribe $f^{2}=f\circ f$, $f^{3}=f\circ f\circ f$. (a) Para $f\colon\Z\to\Z$, $f(n)=2n-1$, calcule $f^{2}$ y $f^{3}$ y conjeture y demuestre por inducción una fórmula para $f^{k}$. (b) Demuestre que si $f^{3}$ es sobreyectiva, $f$ es sobreyectiva. (c) Demuestre que si $f^{3}$ es inyectiva, $f$ es inyectiva.
\fuente{adaptado de Stanford CS103~\cite{cs103ps3}, Problem Set~3, problema~4 «Eventual Bijections».}
\end{ejercicio}

\begin{ejercicio}{3}{eje:4-inv-comp}
Sean $f\colon A\to B$ y $g\colon B\to C$ biyectivas. Demuestre que $(g\circ f)^{-1}=f^{-1}\circ g^{-1}$.
\fuente{adaptado de Hammack~\cite{hammack}, §12.5 «Inverse Functions», ejercicios de la sección.}
\end{ejercicio}

\begin{ejercicio}{3}{eje:4-derecha}
Sea $f\colon A\to B$. Demuestre que $f$ tiene una inversa por la derecha si y solo si es sobreyectiva. ¿Es única la inversa por la derecha? Estudie el caso $f\colon\Z\to\Z$, $f(n)=\lfloor n/2\rfloor$.
\end{ejercicio}

\begin{ejercicio}{3}{eje:4-involucion}
Una función $f\colon A\to A$ es una \emph{involución} si $f\circ f=\id_A$. (a) Demuestre que toda involución es biyectiva y $f^{-1}=f$. (b) Compruebe que son involuciones: $x\mapsto-x$ en $\R$; $x\mapsto1/x$ en $\R\setminus\{0\}$; $X\mapsto A\setminus X$ en $\Pa(A)$; la inversión de cadenas en $\Sigma^{*}$. (c) Dé una biyección de $\{1,2,3\}$ en sí mismo que no sea una involución.
\fuente{adaptado de Hammack~\cite{hammack}, §12.5, ejercicios de la sección.}
\end{ejercicio}

\begin{ejercicio}{3}{eje:4-cancelar}
Sea $f\colon A\to B$. Demuestre que $f$ es sobreyectiva si y solo si para todo conjunto $C$ y todas las funciones $g,h\colon B\to C$, la igualdad $g\circ f=h\circ f$ implica $g=h$. Formule la propiedad análoga para la inyectividad.
\fuente{adaptado de Cambridge, Discrete Mathematics~\cite{camdm}, hoja~8, ejercicio~11.2.2.}
\end{ejercicio}

\begin{ejercicio}{2}{eje:4-capas}
Sean $f_1,f_2\colon\R^{2}\to\R^{2}$, $f_1(x,y)=(x+y,\ x-y)$ y $f_2(u,v)=(2u,\ v+1)$, dos «capas» elementales. (a) Calcule $F=f_2\circ f_1$ y $F(1,2)$. (b) Calcule $f_1^{-1}$ y $f_2^{-1}$. (c) Calcule $F^{-1}$ con la \cref{prop:t4-inv-comp} y compruebe que $F^{-1}(F(1,2))=(1,2)$.
\end{ejercicio}

\begin{ejercicio}{3}{eje:4-tokenizador}
Sea $(\tau,\delta)$ un tokenizador: $\tau\colon\Sigma^{*}\to\V^{*}$, $\delta\colon\V^{*}\to\Sigma^{*}$ con $\delta\circ\tau=\id_{\Sigma^{*}}$. Sea $p=\tau\circ\delta\colon\V^{*}\to\V^{*}$. (a) Demuestre que $p\circ p=p$. (b) Demuestre que $p=\id_{\V^{*}}$ si y solo si $\delta$ es inyectiva. (c) Con $\Sigma=\{\cad{a},\cad{b}\}$, $\V=\{\cad{a},\cad{b},\cad{ab}\}$, $\delta$ la concatenación y $\tau$ la segmentación que toma siempre el token más largo posible de izquierda a derecha, calcule $p((\cad{a},\cad{b}))$.
\end{ejercicio}
